\section{AHSME 1989}
        \begin{enumerate}
        	\item (\textit{AHSME 1989}) $(-1)^{5^{2}}+1^{2^{5}}=$


 $\textrm{(A)}\ -7\qquad\textrm{(B)}\ -2\qquad\textrm{(C)}\ 0\qquad\textrm{(D)}\ 1\qquad\textrm{(E)}\ 57$


 

	\item (\textit{AHSME 1989}) $\sqrt{\frac{1}{9}+\frac{1}{16}}=$


 $\textrm{(A)}\ \frac{1}5\qquad\textrm{(B)}\ \frac{1}4\qquad\textrm{(C)}\ \frac{2}7\qquad\textrm{(D)}\ \frac{5}{12}\qquad\textrm{(E)}\ \frac{7}{12}$


 

	\item (\textit{AHSME 1989}) A square is cut into three rectangles along two lines parallel to a side, as shown. If the perimeter of each of the three rectangles is $24$, then the area of the original square is


 
                    \begin{figure}[H]
                        \centering
                        \includegraphics[width=0.5\textwidth]{Problem/AHSME_I_1989_3_Problem_Figure_1.png}
                    \end{figure}
                    


 $\textrm{(A)}\ 24\qquad\textrm{(B)}\ 36\qquad\textrm{(C)}\ 64\qquad\textrm{(D)}\ 81\qquad\textrm{(E)}\ 96$


 

	\item (\textit{AHSME 1989}) In the figure, $ABCD$ is an isosceles trapezoid with side lengths $AD=BC=5$, $AB=4$, and $DC=10$. The point $C$ is on $\overline{DF}$ and $B$ is the midpoint of hypotenuse $\overline{DE}$ in right triangle $DEF$. Then $CF=$


 
                    \begin{figure}[H]
                        \centering
                        \includegraphics[width=0.5\textwidth]{Problem/AHSME_I_1989_4_Problem_Figure_1.png}
                    \end{figure}
                    


 $\textrm{(A)}\ 3.25\qquad\textrm{(B)}\ 3.5\qquad\textrm{(C)}\ 3.75\qquad\textrm{(D)}\ 4.0\qquad\textrm{(E)}\ 4.25$


 

	\item (\textit{AHSME 1989}) Toothpicks of equal length are used to build a rectangular grid as shown. If the grid is 20 toothpicks high and 10 toothpicks wide, then the number of toothpicks used is


 
                    \begin{figure}[H]
                        \centering
                        \includegraphics[width=0.5\textwidth]{Problem/AHSME_I_1989_5_Problem_Figure_1.png}
                    \end{figure}
                    


 $\textrm{(A)}\ 30\qquad\textrm{(B)}\ 200\qquad\textrm{(C)}\ 410\qquad\textrm{(D)}\ 420\qquad\textrm{(E)}\ 430$


 

	\item (\textit{AHSME 1989}) If $a,\,b>0$ and the triangle in the first quadrant bounded by the coordinate axes and the graph of $ax+by=6$ has area $6$, then $ab=$


 $\textrm{(A)}\ 3\qquad\textrm{(B)}\ 6\qquad\textrm{(C)}\ 12\qquad\textrm{(D)}\ 108\qquad\textrm{(E)}\ 432$


 

	\item (\textit{AHSME 1989}) In $\triangle ABC$, $\angle A = 100^\circ$, $\angle B = 50^\circ$, $\angle C = 30^\circ$, $\overline{AH}$ is an altitude, and $\overline{BM}$ is a median. Then $\angle MHC=$


 
                    \begin{figure}[H]
                        \centering
                        \includegraphics[width=0.5\textwidth]{Problem/AHSME_I_1989_7_Problem_Figure_1.png}
                    \end{figure}
                    


 $\textrm{(A)}\ 15^\circ\qquad\textrm{(B)}\ 22.5^\circ\qquad\textrm{(C)}\ 30^\circ\qquad\textrm{(D)}\ 40^\circ\qquad\textrm{(E)}\ 45^\circ$


 

	\item (\textit{AHSME 1989}) For how many integers $n$ between $1$ and $100$ does $x^{2}+x-n$ factor into the product of two linear factors with integer coefficients?


 $\textrm{(A)}\ 0\qquad\textrm{(B)}\ 1\qquad\textrm{(C)}\ 2\qquad\textrm{(D)}\ 9\qquad\textrm{(E)}\ 10$


 

	\item (\textit{AHSME 1989}) Mr. and Mrs. Zeta want to name baby Zeta so that its monogram (first, middle, and last initials) will be in alphabetical order with no letter repeated. How many such monograms are possible? 


 $\textrm{(A)}\ 276\qquad\textrm{(B)}\ 300\qquad\textrm{(C)}\ 552\qquad\textrm{(D)}\ 600\qquad\textrm{(E)}\ 15600$


 

	\item (\textit{AHSME 1989}) Consider the sequence defined recursively by $u_{1}= a$ (any positive integer), and $u_{n+1}=\frac{-1}{u_{n}+1}$,  $n = 1,2,3,\cdots$. For which of the following values of $n$ must $u_{n}=a$?


 $\textrm{(A)}\ 14\qquad\textrm{(B)}\ 15\qquad\textrm{(C)}\ 16\qquad\textrm{(D)}\ 17\qquad\textrm{(E)}\ 18$


 

	\item (\textit{AHSME 1989}) Let $a$, $b$, $c$, and $d$ be positive integers with $a < 2b$, $b < 3c$, and $c<4d$. If $d<100$, the largest possible value for $a$ is


 $\textrm{(A)}\ 2367\qquad\textrm{(B)}\ 2375\qquad\textrm{(C)}\ 2391\qquad\textrm{(D)}\ 2399\qquad\textrm{(E)}\ 2400$


 

	\item (\textit{AHSME 1989}) The traffic on a certain east-west highway moves at a constant speed of 60 miles per hour in both directions. An eastbound driver passes 20 west-bound vehicles in a five-minute interval. Assume vehicles in the westbound lane are equally spaced. Which of the following is closest to the number of westbound vehicles present in a 100-mile section of highway?


 $\textrm{(A)}\ 100\qquad\textrm{(B)}\ 120\qquad\textrm{(C)}\ 200\qquad\textrm{(D)}\ 240\qquad\textrm{(E)}\ 400$


 

	\item (\textit{AHSME 1989}) Two strips of width 1 overlap at an angle of $\alpha$ as shown. The area of the overlap (shown shaded) is


 
                    \begin{figure}[H]
                        \centering
                        \includegraphics[width=0.5\textwidth]{Problem/AHSME_I_1989_13_Problem_Figure_1.png}
                    \end{figure}
                    


 $\textrm{(A)}\ \sin\alpha\qquad\textrm{(B)}\ \frac{1}{\sin\alpha}\qquad\textrm{(C)}\ \frac{1}{1-\cos\alpha}\qquad\textrm{(D)}\ \frac{1}{\sin^{2}\alpha}\qquad\textrm{(E)}\ \frac{1}{(1-\cos\alpha)^{2}}$


 

	\item (\textit{AHSME 1989}) $\cot 10+\tan 5 =$


 $\textrm{(A)}\ \csc 5\qquad\textrm{(B)}\ \csc 10\qquad\textrm{(C)}\ \sec 5\qquad\textrm{(D)}\ \sec 10\qquad\textrm{(E)}\ \sin 15$


 

	\item (\textit{AHSME 1989}) In $\triangle ABC$, $AB=5$, $BC=7$, $AC=9$, and $D$ is on $\overline{AC}$ with $BD=5$. Find the ratio of $AD:DC$. 


 
                    \begin{figure}[H]
                        \centering
                        \includegraphics[width=0.5\textwidth]{Problem/AHSME_I_1989_15_Problem_Figure_1.png}
                    \end{figure}
                    


 $\textrm{(A)}\ 4:3\qquad\textrm{(B)}\ 7:5\qquad\textrm{(C)}\ 11:6\qquad\textrm{(D)}\ 13:5\qquad\textrm{(E)}\ 19:8$


 

	\item (\textit{AHSME 1989}) A lattice point is a point in the plane with integer coordinates. How many lattice points are on the line segment whose endpoints are (3,17) and (48,281)? (Include both endpoints of the segment in your count.)


 $\textrm{(A)}\ 2\qquad\textrm{(B)}\ 4\qquad\textrm{(C)}\ 6\qquad\textrm{(D)}\ 16\qquad\textrm{(E)}\ 46$


 

	\item (\textit{AHSME 1989}) The perimeter of an equilateral triangle exceeds the perimeter of a square by $1989$ cm. The length of each side of the triangle exceeds the length of each side of the square by $d$ cm. The square has perimeter greater than 0. How many positive integers are NOT a possible value for $d$?


 $\textrm{(A)}\ 0\qquad\textrm{(B)}\ 9\qquad\textrm{(C)}\ 221\qquad\textrm{(D)}\ 663\qquad\textrm{(E)}\ \text{infinitely many}$


 

	\item (\textit{AHSME 1989}) The set of all numbers x for which $x+\sqrt{x^{2}+1}-\frac{1}{x+\sqrt{x^{2}+1}}$ is a rational number is the set of all:


 $\textrm{(A)}\ \text{ integers }x\qquad\textrm{(B)}\ \text{ rational }x\qquad\textrm{(C)}\ \text{ real }x\qquad\textrm{(D)}\ x\text{ for which }\sqrt{x^{2}+1}\text{ is rational}\qquad\textrm{(E)}\ x\text{ for which }x+\sqrt{x^{2}+1}\text{ is rational }$


 

	\item (\textit{AHSME 1989}) A triangle is inscribed in a circle. The vertices of the triangle divide the circle into three arcs of lengths $3$, $4$, and $5$. What is the area of the triangle?


 $\textrm{(A)}\ 6\qquad\textrm{(B)}\ \frac{18}{\pi^{2}}\qquad\textrm{(C)}\ \frac{9}{\pi^{2}}\left(\sqrt{3}-1\right)\qquad\textrm{(D)}\ \frac{9}{\pi^{2}}\left(\sqrt{3}+1\right)\qquad\textrm{(E)}\ \frac{9}{\pi^{2}}\left(\sqrt{3}+3\right)$


 

	\item (\textit{AHSME 1989}) Let $x$ be a real number selected uniformly at random between 100 and 200. If $\lfloor {\sqrt{x}} \rfloor = 12$, find the probability that $\lfloor {\sqrt{100x}} \rfloor = 120$. ($\lfloor {v} \rfloor$ means the greatest integer less than or equal to $v$.)


 $\text{(A)} \ \frac{2}{25} \qquad \text{(B)} \ \frac{241}{2500} \qquad \text{(C)} \ \frac{1}{10} \qquad \text{(D)} \ \frac{96}{625} \qquad \text{(E)} \ 1$


 

	\item (\textit{AHSME 1989}) A square flag has a red cross of uniform width with a blue square in the center on a white background as shown. (The cross is symmetric with respect to each of the diagonals of the square.) If the entire cross (both the red arms and the blue center) takes up 36\% of the area of the flag, what percent of the area of the flag is blue?


 
                    \begin{figure}[H]
                        \centering
                        \includegraphics[width=0.5\textwidth]{Problem/AHSME_I_1989_21_Problem_Figure_1.png}
                    \end{figure}
                    


 $\text{(A)}\ .5\qquad\text{(B)}\ 1\qquad\text{(C)}\ 2\qquad\text{(D)}\ 3\qquad\text{(E)}\ 6$


 

	\item (\textit{AHSME 1989}) A child has a set of 96 distinct blocks. Each block is one of 2 materials (plastic, wood), 3 sizes (small, medium, large), 4 colors (blue, green, red, yellow), and 4 shapes (circle, hexagon, square, triangle). How many blocks in the set are different from the 'plastic medium red circle' in exactly 2 ways? (The 'wood medium red square' is such a block)


 (A) 29 (B) 39 (C) 48 (D) 56 (E) 62


 

	\item (\textit{AHSME 1989}) A particle moves through the first quadrant as follows. During the first minute it moves from the origin to $(1,0)$. Thereafter, it continues to follow the directions indicated in the figure, going back and forth between the positive x and y axes, moving one unit of distance parallel to an axis in each minute. At which point will the particle be after exactly 1989 minutes?


 
                    \begin{figure}[H]
                        \centering
                        \includegraphics[width=0.5\textwidth]{Problem/AHSME_I_1989_23_Problem_Figure_1.png}
                    \end{figure}
                    


 $\text{(A)}\ (35,44)\qquad\text{(B)}\ (36,45)\qquad\text{(C)}\ (37,45)\qquad\text{(D)}\ (44,35)\qquad\text{(E)}\ (45,36)$


 

	\item (\textit{AHSME 1989}) Five people are sitting at a round table. Let $f\geq 0$ be the number of people sitting next to at least 1 female and $m\geq0$ be the number of people sitting next to at least one male. The number of possible ordered pairs $(f,m)$ is


 $\mathrm{(A) \ 7 } \qquad \mathrm{(B) \ 8 } \qquad \mathrm{(C) \ 9 } \qquad \mathrm{(D) \ 10 } \qquad \mathrm{(E) \ 11 }$


 

	\item (\textit{AHSME 1989}) In a certain cross country meet between 2 teams of 5 runners each, a runner who finishes in the $n$th position contributes $n$ to his team's score. The team with the lower score wins. If there are no ties among the runners, how many different winning scores are possible?


 (A) 10 (B) 13 (C) 27 (D) 120 (E) 126


 

	\item (\textit{AHSME 1989}) A regular octahedron is formed by joining the centers of adjoining faces of a cube. The ratio of the volume of the octahedron to the volume of the cube is


 $\mathrm{(A) \frac{\sqrt{3}}{12} } \qquad \mathrm{(B) \frac{\sqrt{6}}{16} } \qquad \mathrm{(C) \frac{1}{6} } \qquad \mathrm{(D) \frac{\sqrt{2}}{8} } \qquad \mathrm{(E) \frac{1}{4} }$


 

	\item (\textit{AHSME 1989}) Let $n$ be a positive integer. If the equation $2x+2y+z=n$ has 28 solutions in positive integers $x$, $y$, and $z$, then $n$ must be either 


 $\mathrm{(A)}\ 14 \text{ or } 15\ \mathrm{(B)}\ 15 \text{ or } 16\ \mathrm{(C)}\ 16 \text{ or } 17\ \mathrm{(D)}\ 17 \text{ or } 18\ \mathrm{(E)}\ 18 \text{ or } 19$


 

	\item (\textit{AHSME 1989}) Find the sum of the roots of $\tan^2x-9\tan x+1=0$ that are between $x=0$ and $x=2\pi$ radians.


 $\mathrm{(A)  \frac{\pi}{2} } \qquad \mathrm{(B) \pi } \qquad \mathrm{(C) \frac{3\pi}{2} } \qquad \mathrm{(D) 3\pi } \qquad \mathrm{(E) 4\pi }$


 

	\item (\textit{AHSME 1989}) Find $\sum_{k=0}^{49}(-1)^k\binom{99}{2k}$, where $\binom{n}{j}=\frac{n!}{j!(n-j)!}$


 (A) $-2^{50}$ (B) $-2^{49}$ (C) 0 (D) $2^{49}$ (E) $2^{50}$


 



 

	\item (\textit{AHSME 1989}) Suppose that 7 boys and 13 girls line up in a row. Let $S$ be the number of places in the row where a boy and a girl are standing next to each other. For example, for the row $GBBGGGBGBGGGBGBGGBGG$ we have that $S=12$. The average value of $S$ (if all possible orders of these 20 people are considered) is closest to


 $\mathrm{(A) 9 } \qquad \mathrm{(B) 10 } \qquad \mathrm{(C) 11 } \qquad \mathrm{(D) 12 } \qquad \mathrm{(E) 13 }$


 

\end{enumerate}